% ==========================================
% SECTION 4: TASK 3 - MANIPULATION
% ==========================================
\section{Task 3: Manipulation}

\begin{frame}{Task 3: Manipulation Coordination}

    The \texttt{coordinator\_task3} node orchestrates the sequence, bridging the data gathered in Task 2 with the actuation logic.

    \vspace{0.1cm}

    \begin{block}{Execution Workflow}
        \begin{enumerate}
            \item \textbf{Data Ingestion:} Parses \texttt{marker\_locations.csv} to retrieve the exact pose of the targets found during exploration.
            \item \textbf{State Sequence:} Enforces the required order of operations:
            {
                \small
                $$ \text{Localize} \xrightarrow{\text{Navigate to Pick}} \text{Pick Cube ID 63} \xrightarrow{\text{Navigate to Place}} \text{Place Cube ID 63} $$
                $$ \text{Next Target} \xrightarrow{\text{Navigate to Pick}} \text{Pick Cube ID 582} \xrightarrow{\text{Navigate to Place}} \text{Place Cube ID 582} $$
            }
            \item \textbf{Visual Feedback:} Publishes dynamic TF frames (\texttt{target\_cube\_frame}) to visualize the manipulation target in RViz in real-time.
        \end{enumerate}
    \end{block}

\end{frame}

\begin{frame}{Task 3: Grasping \& Manipulation Strategy}

    The manipulation pipeline implements a State Machine handling
    \begin{itemize}
        \item \textbf{Grasp Procedure:} \textit{Pre-Grasp} $\rightarrow$ \textit{Grasp} $\rightarrow$ \textit{Post-Grasp}.
        \item \textbf{Release Procedure:} \textit{Pre-Release} $\rightarrow$ \textit{Release} $\rightarrow$ \textit{Post-Release}.
    \end{itemize}

    Between this two main phases the robot perform navigation (handled by the coordinator and the navigation stack). To ensure a \textbf{safe} navigation the manipulator also perform additional phases.
    \begin{itemize}
        \item \textbf{Safe Pose:} Before navigating, the arm moves to a predefined "safe" configuration to avoid collisions with the environment during navigation.
        \item \textbf{Idle:} After moving to the safe pose, the arm enters an idle state, waiting for the coordinator to signal the next command.
    \end{itemize}
    \vspace{0.2cm}


\end{frame}

\begin{frame}{Task 3: Grasp \& Release Selection Pose}
    \begin{block}{Pose Calculation}
        The Grasp Pose is computed by extracting the marker's transform. The choice of the grasping point has been made to ensure a stable grip on the cube, considering collinear and antipodal grasp for the 7cm cube.
    \end{block}

        \begin{columns}
        \begin{column}{0.6\textwidth}
            \begin{exampleblock}{Pose Computation (KDL)}
                Given the marker pose $T_{base}^{marker}$, we compute the approach vector by applying a geometrical offset:
                $$ T_{grasp} = T_{base}^{marker} \cdot T_{rotY(90^\circ)} \cdot T_{offset(-0.4m)} $$
                This aligns the gripper normal with the cube's face.
            \end{exampleblock}
        \end{column}

        \begin{column}{0.4\textwidth}
            \begin{figure}
                \centering
                % Immagine di RViz con le terne cartesiane (Assi colorati) sul cubo
                \includegraphics[width=0.8\linewidth, height=3cm, keepaspectratio]{placeholder.png}
                \caption{TF Frames visualization}
            \end{figure}
        \end{column}
    \end{columns}

    \vspace{0.3cm}
    \textbf{Placement Logic:}
    The release pose is derived from the grasp pose but lifted vertically to ensure the object rests on the table:
    $$ P_{release}.z = P_{grasp}.z + 0.15m $$

\end{frame}

\begin{frame}{Task 3: Precise Placement Pose}
    The placement pose is chosen to be the same as the grasp pose to ensure that the cube is placed in approximately the same position as the one it was picked from. This is because the robot has to place the cube in a specific area of the table, and using the same pose for both picking and placing allows for a more consistent and accurate placement.


    \begin{alertblock}{Precise Approach Logic}
        To compensate for navigation errors, the node reads the real-time distance from the table. If the robot is wrongly positioned ($>1cm$ error), it triggers a closed-loop \textbf{Linear Approach/Backup} maneuver before attempting the grasp.
    \end{alertblock}
\end{frame}